Scientific agents built on large language models are increasingly deployed across real research workflows, yet remain evaluated at the level of whole tasks: a pass rate tells us \emph{whether} an agent succeeds but not \emph{where} in its execution frontier capability is genuinely required. We introduce step-level counterfactual replacement --- surgically substituting the policy at a single step of a frozen agent trajectory while holding context, tool state and remaining budget fixed --- and apply it to scientific-computing agents executing multi-step data-analysis tasks. Across roughly 1,000 controlled runs spanning four replacement policies (frontier resampling, two API-tier models and open-weights models), replacement damage concentrates at a minority of steps: relative to matched control steps, cheaper models lose 0.21--0.71 success probability at identified signal steps, and the damage is model-specific rather than graded, splitting into protocol lapses, semantic errors and model-idiosyncratic error families. Full-step scans separate healthy, capability-dependent and cliff steps; structured feedback partially repairs failures (+0.21 success rate) whereas naive step-splicing amplifies them, and doubling model scale does not dissolve the failures. Finally, two annotators blind to all outcomes code verifiability features of 74 step instances from pre-registered materials; these features predict step-level replaceability under task-level cross-validation, meeting both pre-registered criteria on the original 38 points ($\Delta$$\rho$ +0.17) and reproducing them on the enlarged 49-point set ($\Delta$$\rho$ +0.13, paired permutation p = 0.001); the gain, however, is an in-distribution effect --- it reverses on unseen instances of the training tasks and turns negative on a fully unseen domain (computational chemistry, 18 points, 72 observations: $\Delta$$\rho$ between $-$0.18 and $-$0.27 across training variants, the feature model consistently worse than the arm-prior control in rank correlation and worse in RMSE under every configuration). Verifiability profiling turns "which model for which step" from anecdote into measurement, and offers a design principle for allocating frontier capability in scientific agents.
